\documentclass[12pt,a4paper]{article}
\usepackage[utf-8]{inputenc}
\usepackage[margin=1in]{geometry}
\usepackage{amsmath}
\usepackage{amssymb}
\usepackage{graphicx}
\usepackage{hyperref}
\usepackage{xcolor}
\usepackage{listings}
\usepackage{setspace}
\usepackage{fancyhdr}

% Set line spacing
\onehalfspacing

% Header and footer
\pagestyle{fancy}
\fancyhf{}
\lhead{MDP Deep Dive: iRDRC Application}
\rhead{Discussion 2}
\cfoot{\thepage}

% Title
\title{\textbf{Markov Decision Process Deep Dive}\\
\large Understanding Sequential Decision-Making in Autonomous Vehicles\\
\normalsize iRDRC Application}
\author{Based on: Hieu, Hoang, Luong, Niyato (IEEE WCL 2020)}
\date{2026}

\begin{document}

\maketitle
\setcounter{tocdepth}{1}
\tableofcontents
\newpage

% ============================================================================
\section{What is a Markov Decision Process?}
% ============================================================================

A \textbf{Markov Decision Process (MDP)} is a mathematical framework for modeling sequential decision-making problems in stochastic (probabilistic) environments. It answers the question: \textbf{How should an agent make optimal decisions over time when the future is uncertain?}

\subsection{Core Motivation}

Many real-world problems require decisions at multiple time steps:
\begin{itemize}
    \item A self-driving car must decide whether to brake, accelerate, or change lanes
    \item An autonomous vehicle must allocate limited hardware (radar vs. communication)
    \item A robot must plan which actions to take to reach a goal
\end{itemize}

In all these cases:
\begin{enumerate}
    \item The current decision affects future state (state transitions)
    \item The future is uncertain (randomness in the environment)
    \item We want to maximize long-term success, not just immediate reward
\end{enumerate}

The MDP framework formalizes this problem.

\subsection{The MDP Tuple: $\langle \mathcal{S}, \mathcal{A}, R, P, \gamma \rangle$}

An MDP is completely defined by five components:

\[
\boxed{\text{MDP} = \langle \mathcal{S}, \mathcal{A}, R, P, \gamma \rangle}
\]

\begin{table}[h]
\centering
\begin{tabular}{|c|l|}
\hline
\textbf{Component} & \textbf{Meaning} \\
\hline
$\mathcal{S}$ & State space: all possible situations the agent can be in \\
\hline
$\mathcal{A}$ & Action space: all possible decisions the agent can make \\
\hline
$R: \mathcal{S} \times \mathcal{A} \to \mathbb{R}$ & Reward function: immediate payoff for each action \\
\hline
$P: \mathcal{S} \times \mathcal{A} \times \mathcal{S} \to [0,1]$ & Transition probabilities: probability of reaching state $s'$ from $s$ after action $a$ \\
\hline
$\gamma \in [0,1)$ & Discount factor: how much we value future rewards \\
\hline
\end{tabular}
\end{table}

\subsection{The Markov Property}

The name ``Markov'' refers to a key assumption: \textbf{the future depends only on the present, not the past}.

Formally, given the current state $s_t$, the next state $s_{t+1}$ is independent of all previous states:

\[
P(s_{t+1} | s_t, a_t, s_{t-1}, a_{t-1}, \ldots, s_0, a_0) = P(s_{t+1} | s_t, a_t)
\]

\textbf{Intuition}: The current state contains all relevant information to predict the future. We don't need to remember history.

\textbf{Example}: If an AV is at position $(x, y)$ with velocity $v$, we can predict its next position without knowing how it got to $(x, y)$.

---

\section{Formulating iRDRC as an MDP}
\label{sec:irdrc_mdp}

Now let's concretely apply the MDP framework to the iRDRC problem. Recall: an autonomous vehicle must decide at each time step whether to operate in \textbf{Radar Mode} (sensing) or \textbf{Communication Mode} (transmitting data).

\subsection{The Problem Context}

The AV faces a tradeoff:
\begin{itemize}
    \item \textbf{Radar Mode}: Detect hazards (safety), but can't transmit data
    \item \textbf{Communication Mode}: Send/receive data, but may miss hazards
\end{itemize}

The MDP will help us find the \textbf{optimal policy}: a decision rule that balances these competing objectives.

\subsection{Component 1: State Space $\mathcal{S}$}

At time $t$, the AV observes six pieces of information:

\[
s_t = (d_t, c_t, r_t, w_t, v_t, m_t)
\]

\begin{table}[h]
\centering
\begin{tabular}{|c|c|l|}
\hline
\textbf{Symbol} & \textbf{Range} & \textbf{Meaning} \\
\hline
$d_t$ & $\{0, 1, \ldots, D\}$ & Queue size (packets waiting to transmit) \\
\hline
$c_t$ & $\{0, 1\}$ & Channel quality: 0=bad, 1=good \\
\hline
$r_t$ & $\{0, 1\}$ & Road condition: 0=good, 1=hazardous/slippery \\
\hline
$w_t$ & $\{0, 1\}$ & Weather: 0=clear, 1=bad (rain, fog, snow) \\
\hline
$v_t$ & $\{0, 1\}$ & AV speed: 0=low ($< 60$ km/h), 1=high ($\geq 60$ km/h) \\
\hline
$m_t$ & $\{0, 1\}$ & Moving object: 0=no, 1=car/pedestrian nearby \\
\hline
\end{tabular}
\end{table}

\subsubsection{State Space Size}

With $D = 10$ (queue capacity), the total number of possible states is:

\[
|\mathcal{S}| = (D+1) \times 2^5 = 11 \times 32 = \boxed{352 \text{ states}}
\]

This is tractable for many learning algorithms.

\subsection{Component 2: Action Space $\mathcal{A}$}

At each time step, the AV chooses exactly one of two actions:

\[
\mathcal{A} = \{0, 1\}
\]

\begin{table}[h]
\centering
\begin{tabular}{|c|l|}
\hline
\textbf{Action} & \textbf{Meaning} \\
\hline
$a_t = 0$ & Communication Mode (transmit data) \\
\hline
$a_t = 1$ & Radar Mode (sense environment) \\
\hline
\end{tabular}
\end{table}

Simple binary choice, but the consequences ripple through time.

\subsection{Component 3: Reward Function $R(s, a)$}

The reward function defines what behavior we want the policy to learn. It's the ``objective function'' of the problem.

The paper designs the reward function to balance two competing goals:
\begin{enumerate}
    \item \textbf{Safety}: Detect unexpected hazards (prioritized)
    \item \textbf{Throughput}: Transmit data packets (secondary)
\end{enumerate}

\subsubsection{Reward Assignment}

The reward depends on:
\begin{itemize}
    \item \textbf{Action taken}: $a_t$ (radar or communication)
    \item \textbf{Channel quality}: $c_t$ (affects communication success)
    \item \textbf{Unexpected event}: Whether $\oplus$ occurs
\end{itemize}

We use notation:
\begin{itemize}
    \item $\oplus$ = unexpected event occurs (hazard on road)
    \item $\bar{\oplus}$ = no unexpected event (safe environment)
\end{itemize}

\textbf{Case 1: Communication Mode, Bad Channel, Event Occurs}
\begin{align}
a_t = 0, \quad c_t = 0, \quad \oplus &\implies r_t = +r_1 = +2 \\
&\text{(Data transmission failed, but tried; missed hazard)}
\end{align}

\textbf{Case 2: Communication Mode, Good Channel, Event Occurs}
\begin{align}
a_t = 0, \quad c_t = 1, \quad \oplus &\implies r_t = +r_2 = +50 \\
&\text{(Data transmitted successfully, but missed hazard)}
\end{align}

\textbf{Case 3: Radar Mode, Event Occurs}
\begin{align}
a_t = 1, \quad \oplus &\implies r_t = -r_3 = -50 \\
&\text{(Detected hazard, but hazard occurred $\Rightarrow$ negative reward)}
\end{align}

\textbf{Case 4: Radar Mode, No Event Occurs}
\begin{align}
a_t = 1, \quad \bar{\oplus} &\implies r_t = +r_4(b+1) \\
&\text{where } b = \text{number of consecutive radar steps} \\
&\text{(Safe environment; proactive sensing rewarded)}
\end{align}

\subsubsection{Reward Function Summary}

\begin{table}[h]
\centering
\begin{tabular}{|c|c|c|c|}
\hline
\textbf{Action} & \textbf{Channel} & \textbf{Event?} & \textbf{Reward} \\
\hline
Comm & Bad & $\oplus$ & $r_1 = 2$ \\
\hline
Comm & Good & $\oplus$ & $r_2 = 50$ \\
\hline
Radar & N/A & $\oplus$ & $-r_3 = -50$ \\
\hline
Radar & N/A & $\bar{\oplus}$ & $r_4(b+1)$ (typically $r_4 = 5$) \\
\hline
\end{tabular}
\end{table}

\textbf{Key Design Insight}: Notice that $r_3 = 50$ (magnitude) is \textbf{much larger} than $r_1 = 2$ or $r_2 = 50$. This heavily penalizes missing a hazard in communication mode, enforcing that \textbf{safety is the priority}.

\subsection{Component 4: State Transition Probabilities $P(s' | s, a)$}

When the AV takes action $a$ in state $s$, the environment stochastically transitions to state $s'$. The transition probability captures this randomness.

\subsubsection{Queue Dynamics}

The data queue evolves as:
\begin{itemize}
    \item If $a_t = 0$ (communication): $d_{t+1} = d_t - 1$ (transmit one packet)
    \item If $a_t = 1$ (radar): $d_{t+1} = d_t + n$ (new packets arrive)
\end{itemize}

This creates a tension: radar improves safety but lets data accumulate.

\subsubsection{Environmental Factor Transitions}

Environmental factors $(r, w, v, m)$ follow a Markov process (their next values depend only on current values, not history). The exact transition probabilities are not specified in the paper, but in practice they are estimated from real-world data.

\subsubsection{Unexpected Event Probability}

The probability that an unexpected event occurs is computed using Bayes' theorem:

\[
P(\oplus \mid r_t, w_t, v_t, m_t) = \sum_{i \in \{r,w,v,m\}} p_i^{s_i(t)}
\]

where $p_i^{s_i}$ is the conditional probability of an unexpected event given factor $i$ is in state $s_i$.

\textbf{Example}: If $p_v^1 = 0.1, p_w^1 = 0.15, p_r^0 = 0, p_m^1 = 0.05$, and current state is $(r=0, w=1, v=1, m=1)$, then:

\[
P(\oplus) = 0 + 0.15 + 0.1 + 0.05 = 0.30
\]

There's a 30\% chance of an unexpected event; the policy should account for this.

\subsection{Component 5: Discount Factor $\gamma$}

The discount factor $\gamma \in [0, 1)$ specifies how much we care about future rewards versus immediate rewards.

\[
\gamma = 0.99 \quad (\text{typical value})
\]

\textbf{Interpretation}: A reward received $n$ steps in the future is worth $\gamma^n$ times as much as an immediate reward.

\[
\text{Value of reward in 10 steps} = \gamma^{10} \times \text{immediate reward} = 0.99^{10} \approx 0.9 \times \text{immediate reward}
\]

This makes the AV prefer immediate safety but not at the cost of completely ignoring long-term performance.

---

\section{Solving the MDP: Policies and Value Functions}
\label{sec:solving}

Now that we've formulated the problem as an MDP, the question is: \textbf{How do we find the best policy?}

\subsection{What is a Policy?}

A \textbf{policy} $\pi$ is a decision rule that maps states to actions:

\[
\pi: \mathcal{S} \to \mathcal{A}
\]

For each possible state $s$, the policy tells us which action to take: $a = \pi(s)$.

\textbf{Example policy} (deterministic):
\begin{align}
\pi(s) = \begin{cases}
1 & \text{if } w_t = 1 \text{ or } r_t = 1 \text{ (bad weather or road $\Rightarrow$ radar)} \\
0 & \text{otherwise (communication)}
\end{cases}
\end{align}

This is a \textbf{heuristic policy}: it uses domain knowledge but may not be optimal.

\subsection{Evaluating a Policy: The Return}

Given a policy $\pi$, how do we measure its performance? We compute the \textbf{discounted return}:

\[
G_t = \sum_{k=0}^{\infty} \gamma^k r_{t+k}
\]

This is the sum of all future rewards, with distant rewards discounted.

For finite horizons ($T$ steps):
\[
G_t = \sum_{k=0}^{T-1} \gamma^k r_{t+k} + \text{(terminal value)}
\]

\subsection{Value Functions}

To compare policies, we need to measure expected performance. This is where \textbf{value functions} come in.

\subsubsection{State Value Function $V^{\pi}(s)$}

The \textbf{state value} is the expected return starting from state $s$ and following policy $\pi$:

\[
V^{\pi}(s) = \mathbb{E}\left[ G_t \mid s_t = s, \pi \right] = \mathbb{E}\left[ \sum_{k=0}^{\infty} \gamma^k r_{t+k} \mid s_t = s, \pi \right]
\]

\textbf{Intuition}: How good is it to be in state $s$ if we follow policy $\pi$?

\textbf{Example}: If $V^{\pi}(s_1) = 100$ and $V^{\pi}(s_2) = 50$, then state $s_1$ is ``better'' under policy $\pi$ (higher expected future reward).

\subsubsection{Action Value Function $Q^{\pi}(s, a)$}

The \textbf{action value} (or Q-value) is the expected return from taking action $a$ in state $s$, then following policy $\pi$:

\[
Q^{\pi}(s, a) = \mathbb{E}\left[ G_t \mid s_t = s, a_t = a, \pi \right]
\]

\textbf{Intuition}: How good is it to take action $a$ in state $s$?

\subsubsection{Relationship Between $V$ and $Q$}

\[
V^{\pi}(s) = \sum_{a \in \mathcal{A}} \pi(a|s) Q^{\pi}(s, a)
\]

The state value is a weighted average of action values, weighted by the policy's probability of taking each action.

\subsection{The Bellman Equations}

These fundamental recursive equations express value functions in terms of immediate rewards and future values.

\subsubsection{Bellman Equation for $V^{\pi}(s)$}

\[
V^{\pi}(s) = \sum_{a \in \mathcal{A}} \pi(a|s) \sum_{s'} P(s'|s,a) \left[ R(s,a) + \gamma V^{\pi}(s') \right]
\]

\textbf{Breakdown}:
\begin{itemize}
    \item $\pi(a|s)$: probability of taking action $a$ in state $s$
    \item $R(s,a)$: immediate reward for action $a$ in state $s$
    \item $\gamma V^{\pi}(s')$: discounted value of future state $s'$
    \item $P(s'|s,a)$: probability of transitioning to $s'$
\end{itemize}

\textbf{Intuition}: The value of a state is the immediate reward, plus the expected value of the next state (discounted).

\subsubsection{Bellman Equation for $Q^{\pi}(s,a)$}

\[
Q^{\pi}(s, a) = R(s, a) + \gamma \sum_{s'} P(s'|s,a) V^{\pi}(s')
\]

\subsection{Optimal Policy and Optimal Value Functions}

The \textbf{optimal policy} $\pi^*$ is the one that maximizes expected return:

\[
\pi^* = \arg\max_{\pi} V^{\pi}(s) \quad \forall s \in \mathcal{S}
\]

The \textbf{optimal value function} is:

\[
V^*(s) = \max_{\pi} V^{\pi}(s) = \mathbb{E}[G_t | s_t = s, \pi^*]
\]

And the optimal action-value function:

\[
Q^*(s, a) = \max_{\pi} Q^{\pi}(s, a)
\]

\subsubsection{Optimal Bellman Equations}

\[
V^*(s) = \max_a \sum_{s'} P(s'|s,a) \left[ R(s,a) + \gamma V^*(s') \right]
\]

\[
Q^*(s, a) = R(s, a) + \gamma \sum_{s'} P(s'|s,a) \max_{a'} Q^*(s', a')
\]

These are the recursive equations that define optimality.

---

\section{Solving MDPs: Algorithms}

Several algorithms can solve MDPs (find $\pi^*$). iRDRC uses \textbf{Deep Q-Learning (DQN)}.

\subsection{Why Not Tabular Q-Learning?}

\textbf{Tabular Q-Learning} maintains a table $Q[s, a]$ for each state-action pair and updates it as:

\[
Q[s,a] \leftarrow Q[s,a] + \alpha \left( r + \gamma \max_{a'} Q[s', a'] - Q[s,a] \right)
\]

With 352 states, this is feasible. But tabular methods fail when:
\begin{itemize}
    \item State space is continuous or very large
    \item Need to generalize to unseen states
\end{itemize}

\subsection{Deep Q-Networks (DQN) for iRDRC}

DQN uses a neural network to approximate Q-values. Instead of maintaining a table, we have:

\[
Q(s, a) \approx Q_\theta(s, a)
\]

where $\theta$ are the network weights.

\textbf{Training procedure}:
\begin{enumerate}
    \item Agent takes action, observes reward and next state
    \item Store transition $(s, a, r, s')$ in replay buffer
    \item Sample mini-batch from buffer
    \item Compute target: $y = r + \gamma \max_{a'} Q_{\theta'}(s', a')$
    \item Update weights to minimize loss: $L = (y - Q_\theta(s, a))^2$
\end{enumerate}

Key techniques:
\begin{itemize}
    \item \textbf{Experience Replay}: Sample past transitions randomly (breaks correlation)
    \item \textbf{Target Network}: Separate network for computing targets (stabilizes training)
\end{itemize}

---

\section{Summary: MDP Framework for iRDRC}

\begin{table}[h]
\centering
\begin{tabular}{|l|l|}
\hline
\textbf{Component} & \textbf{iRDRC Definition} \\
\hline
\textbf{States} & $(d, c, r, w, v, m)$ with 352 total states \\
\hline
\textbf{Actions} & Radar (1) or Communication (0) \\
\hline
\textbf{Rewards} & $r_1=2, r_2=50, -r_3=-50, r_4(b+1)$ with $r_4=5$ \\
\hline
\textbf{Transitions} & Queue dynamics + environmental Markov process \\
\hline
\textbf{Discount} & $\gamma = 0.99$ \\
\hline
\textbf{Goal} & Find policy $\pi^*$ that maximizes $\mathbb{E}[G_t]$ \\
\hline
\textbf{Solution} & Deep Q-Network with experience replay \\
\hline
\end{tabular}
\end{table}

---

\appendix

\section{Quick Reference: Notation}

\begin{tabular}{|c|l|}
\hline
$s_t$ & State at time $t$ \\
\hline
$a_t$ & Action at time $t$ \\
\hline
$r_t$ & Immediate reward at time $t$ \\
\hline
$G_t$ & Discounted return (cumulative reward) \\
\hline
$\gamma$ & Discount factor (0.99) \\
\hline
$\pi$ & Policy (state $\to$ action mapping) \\
\hline
$\pi^*$ & Optimal policy \\
\hline
$V(s)$ & State value (expected return from $s$) \\
\hline
$Q(s,a)$ & Action value (expected return from $(s,a)$) \\
\hline
$P(s'|s,a)$ & Transition probability \\
\hline
$\oplus$ & Unexpected event occurs \\
\hline
$\bar{\oplus}$ & No unexpected event \\
\hline
\end{tabular}

\end{document}
